\section{微积分综合复习}

\subsection{本章速览}

\begin{itemize}
\item \textbf{极限与连续：} 等价无穷小只能替换乘除因子；洛必达法则要求 $\frac00$ 或 $\frac{\infty}{\infty}$ 型且导数之比的极限存在（或为无穷）；用泰勒公式时，展开阶数要与分母的最低阶相匹配。
\item \textbf{一元微分与微分方程：} 可导必连续，连续不一定可导；隐函数求导是对恒等式两边同时求导；极值先用必要条件 $f'(x)=0$ 或 $f'(x)$ 不存在筛出候选点，再用一阶变号或二阶导数判别；一阶线性方程 $y'+P(x)y=Q(x)$ 用积分因子 $e^{\int P\,\mathrm{d}x}$ 求解，二阶常系数非齐次方程按特征根与右端类型设特解，共振时乘以 $x^k$。
\item \textbf{一元积分：} 变上限积分求导公式为 $\dfrac{\mathrm{d}}{\mathrm{d}x}\displaystyle\int_a^{\varphi(x)}f(t)\,\mathrm{d}t=f(\varphi(x))\varphi'(x)$；反常积分要在每个瑕点和无穷远分别判敛，常用标杆 $\displaystyle\int_1^{+\infty}\frac{\mathrm{d}x}{x^p}$（$p>1$ 收敛）与 $\displaystyle\int_0^1\frac{\mathrm{d}x}{x^p}$（$p<1$ 收敛）。
\item \textbf{中值定理：} 罗尔定理要求闭区间连续、开区间可导、两端函数值相等；拉格朗日定理为 $f(b)-f(a)=f'(\xi)(b-a)$；泰勒公式带拉格朗日余项时要求函数有直到 $n+1$ 阶的导数，估计不等式常从余项入手。
\item \textbf{级数：} 正项级数用比较、比值、根值判别；交错级数用莱布尼茨判别，要求 $u_n$ 单调递减且趋于零；幂级数先求收敛半径，再单独检查端点，和函数可在收敛区间内逐项求导或积分，端点值用阿贝尔定理。
\item \textbf{多元微积分：} 偏导数存在不保证可微，可微时才有全微分 $\mathrm{d}z=f_x\,\mathrm{d}x+f_y\,\mathrm{d}y$；二重积分换序要重画区域，极坐标下 $\mathrm{d}x\mathrm{d}y=r\,\mathrm{d}r\mathrm{d}\theta$；格林、高斯、斯托克斯公式都要求偏导连续、曲线或曲面定向正确，高斯公式中的散度为 $P_x+Q_y+R_z$。
\end{itemize}

\subsection{基础过关题}

\textbf{1.} 求极限
\[
\lim_{x\to 0}\left(\frac{1}{\ln(1+x)}-\frac{1}{x}\right).
\]

\ifshowsolutions%
\boxed{\textbf{解}}

\textbf{思路：} 先通分化为 $\frac00$ 型，再把分子、分母都展开到 $x^2$ 项，最低阶相同的项决定了极限。

通分得
\[
\frac{1}{\ln(1+x)}-\frac{1}{x}
=\frac{x-\ln(1+x)}{x\ln(1+x)}.
\]
由带佩亚诺余项的泰勒公式，
\[
\ln(1+x)=x-\frac{x^2}{2}+\frac{x^3}{3}+o(x^3),
\]
故
\[
x-\ln(1+x)=\frac{x^2}{2}-\frac{x^3}{3}+o(x^3),
\]
\[
x\ln(1+x)=x\left(x-\frac{x^2}{2}+o(x^2)\right)
=x^2-\frac{x^3}{2}+o(x^3).
\]
于是
\[
\frac{x-\ln(1+x)}{x\ln(1+x)}
=\frac{\frac{x^2}{2}-\frac{x^3}{3}+o(x^3)}
{x^2-\frac{x^3}{2}+o(x^3)}
=\frac{\frac12-\frac{x}{3}+o(x)}{1-\frac{x}{2}+o(x)}
\longrightarrow \frac12.
\]
\boxed{\text{答案：}\dfrac12}
\else \answerarea[5] \fi%

\textbf{2.} 设 $y=y(x)$ 由方程
\[
e^y+xy=e
\]
确定，求 $y''(0)$。

\ifshowsolutions%
\boxed{\textbf{解}}

\textbf{思路：} 先由原方程求出 $y(0)$，再对恒等式两边连续求导两次，每求一次就代入 $x=0$ 的数值，避免把 $y'$ 的表达式越写越复杂。

当 $x=0$ 时，$e^{y(0)}=e$，故 $y(0)=1$。

对 $e^y+xy=e$ 两边关于 $x$ 求导：
\[
e^y y'+y+xy'=0. \tag{1}
\]
代入 $x=0,\ y=1$，得
\[
e\,y'(0)+1=0,\qquad y'(0)=-\frac1e.
\]
再对 (1) 两边求导：
\[
e^y(y')^2+e^y y''+y'+y'+xy''=0,
\]
即
\[
(e^y+x)y''+e^y(y')^2+2y'=0. \tag{2}
\]
代入 $x=0,\ y=1,\ y'(0)=-\frac1e$：
\[
e\,y''(0)+e\cdot\frac{1}{e^2}-\frac{2}{e}=0,
\]
即
\[
e\,y''(0)=\frac{1}{e},\qquad y''(0)=\frac{1}{e^2}.
\]
\boxed{\text{答案：}y''(0)=e^{-2}}
\else \answerarea[5] \fi%

\textbf{3.} 计算定积分
\[
\int_0^1 x\arctan x\,\mathrm{d}x.
\]

\ifshowsolutions%
\boxed{\textbf{解}}

\textbf{思路：} 被积函数是“多项式乘反三角函数”，用分部积分，把 $\arctan x$ 留下求导、把 $x\,\mathrm{d}x$ 凑成 $\mathrm{d}(x^2/2)$。

取 $u=\arctan x$，$\mathrm{d}v=x\,\mathrm{d}x$，则 $v=\dfrac{x^2}{2}$。分部积分得
\[
\int_0^1 x\arctan x\,\mathrm{d}x
=\left[\frac{x^2}{2}\arctan x\right]_0^1
-\frac12\int_0^1\frac{x^2}{1+x^2}\,\mathrm{d}x.
\]
第一项为 $\dfrac12\cdot\dfrac{\pi}{4}=\dfrac{\pi}{8}$。又
\[
\frac{x^2}{1+x^2}=1-\frac{1}{1+x^2},
\]
故
\[
\int_0^1\frac{x^2}{1+x^2}\,\mathrm{d}x
=\int_0^1\left(1-\frac{1}{1+x^2}\right)\mathrm{d}x
=1-\frac{\pi}{4}.
\]
代回得
\[
\int_0^1 x\arctan x\,\mathrm{d}x
=\frac{\pi}{8}-\frac12\left(1-\frac{\pi}{4}\right)
=\frac{\pi}{4}-\frac12.
\]
\boxed{\text{答案：}\dfrac{\pi}{4}-\dfrac12}
\else \answerarea[6] \fi%

\textbf{4.} 求幂级数
\[
\sum_{n=1}^{\infty}\frac{(x-1)^n}{n\,3^n}
\]
的收敛域与和函数。

\ifshowsolutions%
\boxed{\textbf{解}}

\textbf{思路：} 令 $t=\dfrac{x-1}{3}$ 化为标准形 $\sum \frac{t^n}{n}$，先求关于 $t$ 的收敛区间，再换回 $x$，最后单独检查两个端点。

令 $t=\dfrac{x-1}{3}$，级数变为 $\displaystyle\sum_{n=1}^{\infty}\frac{t^n}{n}$。由比值判别法，
\[
\lim_{n\to\infty}\left|\frac{t^{n+1}/(n+1)}{t^n/n}\right|
=|t|\lim_{n\to\infty}\frac{n}{n+1}=|t|,
\]
故当 $|t|<1$ 时绝对收敛，当 $|t|>1$ 时发散。

检查端点：$t=1$ 时级数为调和级数 $\sum \frac1n$，发散；$t=-1$ 时为交错级数 $\sum\frac{(-1)^n}{n}$，由莱布尼茨判别法收敛。因此关于 $t$ 的收敛域为 $[-1,1)$，换回 $x$ 得
\[
-1\le \frac{x-1}{3}<1
\quad\Longleftrightarrow\quad
-2\le x<4.
\]

在 $|t|<1$ 时，由 $\displaystyle\sum_{n=1}^{\infty}\frac{t^n}{n}=-\ln(1-t)$，得
\[
S(x)=-\ln\left(1-\frac{x-1}{3}\right)
=\ln\frac{3}{4-x},\qquad -2\le x<4.
\]
当 $x=-2$ 即 $t=-1$ 时，上式给出 $\ln\dfrac12=-\ln2$，而级数为 $\sum\frac{(-1)^n}{n}=-\ln2$，由阿贝尔定理和函数的连续性，端点值仍可由该式给出。
\boxed{\text{答案：}\text{收敛域 }[-2,4),\ S(x)=\ln\frac{3}{4-x}}
\else \answerarea[8] \fi%

\textbf{5.} 设曲面 $\Sigma$ 由方程 $z=x^2+xy+y^2$ 确定，点 $P(1,1,3)$ 在 $\Sigma$ 上。求 $\Sigma$ 在 $P$ 处的切平面方程与法线方程。

\ifshowsolutions%
\boxed{\textbf{解}}

\textbf{思路：} 把曲面写成隐式 $F(x,y,z)=0$，则法向量为 $\nabla F(P)$；也可把 $z$ 看成 $x,y$ 的函数，法向量取 $(z_x,z_y,-1)$。

令
\[
F(x,y,z)=x^2+xy+y^2-z.
\]
则
\[
F_x=2x+y,\qquad F_y=x+2y,\qquad F_z=-1.
\]
在 $P(1,1,3)$ 处，
\[
\nabla F(P)=(3,\,3,\,-1).
\]
切平面方程为
\[
3(x-1)+3(y-1)-(z-3)=0,
\]
整理得
\[
3x+3y-z-3=0.
\]
法线方程为
\[
\frac{x-1}{3}=\frac{y-1}{3}=\frac{z-3}{-1}.
\]
\boxed{\text{答案：}3x+3y-z-3=0,\ \frac{x-1}{3}=\frac{y-1}{3}=\frac{z-3}{-1}}
\else \answerarea[6] \fi%

\subsection{综合提高题}

\textbf{6.} 设 $a$ 为实参数，讨论方程
\[
x^3-3x+a=0
\]
的互不相同实根个数。

\ifshowsolutions%
\boxed{\textbf{解}}

\textbf{思路：} 先说明实根至多三个，再用导数确定函数的单调区间与极值，最后按极值的正负讨论每个单调区间内是否穿过 $x$ 轴。

令 $f_a(x)=x^3-3x+a$。若方程有四个互不相同的实根，则由罗尔定理，$f_a'$ 至少有三个互不相同的零点。但
\[
f_a'(x)=3x^2-3=3(x-1)(x+1)
\]
只有 $x=\pm1$ 两个零点，矛盾。因此互不相同的实根至多三个。

由 $f_a'$ 的符号，$f_a$ 在 $(-\infty,-1)$ 上单调递增，在 $(-1,1)$ 上单调递减，在 $(1,+\infty)$ 上单调递增。极大值与极小值分别为
\[
f_a(-1)=2+a,\qquad f_a(1)=a-2.
\]
又
\[
\lim_{x\to-\infty}f_a(x)=-\infty,\qquad
\lim_{x\to+\infty}f_a(x)=+\infty.
\]

分情况讨论：

（1）当 $|a|>2$ 时，极大值与极小值同号。若 $a>2$，则 $f_a(-1)>0$ 且 $f_a(1)>0$，只在 $(-\infty,-1)$ 内有一个根；若 $a<-2$，则 $f_a(-1)<0$ 且 $f_a(1)<0$，只在 $(1,+\infty)$ 内有一个根。此时恰有一个实根。

（2）当 $a=2$ 时，$f_a(1)=0$，而 $f_a(-1)=4>0$，故 $x=1$ 是一个根，另一个根在 $(-\infty,-1)$ 内，共两个实根。当 $a=-2$ 时，$f_a(-1)=0$，而 $f_a(1)=-4<0$，故 $x=-1$ 是一个根，另一个根在 $(1,+\infty)$ 内，也共两个实根。

（3）当 $-2<a<2$ 时，$f_a(-1)=a+2>0$，$f_a(1)=a-2<0$。由零点定理，三个区间 $(-\infty,-1)$、$(-1,1)$、$(1,+\infty)$ 内各至少有一个根；再由每个区间上的严格单调性，各恰有一个根，共三个实根。

\boxed{\text{答案：}|a|>2:1\text{ 个};\ |a|=2:2\text{ 个};\ |a|<2:3\text{ 个}}
\else \answerarea[10] \fi%

\textbf{7.} 设 $f(x)$ 在 $[0,+\infty)$ 上连续，且满足
\[
f(x)=e^x+\int_0^x (x-t)f(t)\,\mathrm{d}t.
\]
求 $f(x)$。

\ifshowsolutions%
\boxed{\textbf{解}}

\textbf{思路：} 含 $(x-t)$ 的积分可拆成 $x\int f-\int tf$，求导一次会消去 $x$ 因子；连续两次求导即可把积分方程化为二阶常系数线性方程。

将积分拆开：
\[
f(x)=e^x+x\int_0^x f(t)\,\mathrm{d}t-\int_0^x tf(t)\,\mathrm{d}t.
\]
由于 $f$ 连续，上式右端可导，故
\[
f'(x)=e^x+\int_0^x f(t)\,\mathrm{d}t+x f(x)-x f(x)
=e^x+\int_0^x f(t)\,\mathrm{d}t.
\]
再求导得
\[
f''(x)=e^x+f(x),
\]
即
\[
f''(x)-f(x)=e^x. \tag{1}
\]
初始条件为
\[
f(0)=e^0+0=1,\qquad f'(0)=e^0+\int_0^0f(t)\,\mathrm{d}t=1.
\]

齐次方程 $f''-f=0$ 的特征根为 $r=\pm1$，齐次通解为
\[
C_1e^x+C_2e^{-x}.
\]
右端 $e^x$ 对应的指数 $1$ 是特征根，故设特解 $f^*(x)=Axe^x$。代入 (1)：
\[
(f^*)''-f^*=(Ax+2A)e^x-Axe^x=2Ae^x=e^x,
\]
得 $A=\dfrac12$。于是
\[
f(x)=C_1e^x+C_2e^{-x}+\frac{x}{2}e^x.
\]
由 $f(0)=1$ 得 $C_1+C_2=1$；由
\[
f'(x)=C_1e^x-C_2e^{-x}+\frac{1+x}{2}e^x
\]
及 $f'(0)=1$ 得 $C_1-C_2+\dfrac12=1$，即 $C_1-C_2=\dfrac12$。解得
\[
C_1=\frac34,\qquad C_2=\frac14.
\]
故
\[
f(x)=\frac34e^x+\frac14e^{-x}+\frac{x}{2}e^x.
\]

最后验证它确实满足原积分方程。令
\[
G(x)=f(x)-e^x-\int_0^x(x-t)f(t)\,\mathrm{d}t.
\]
由上式及初始条件，$G(0)=1-1=0$，$G'(0)=f'(0)-1=0$；并且
\[
G''(x)=f''(x)-e^x-f(x)=0.
\]
故 $G(x)\equiv0$，即所求函数满足原方程。
\boxed{\text{答案：}f(x)=\dfrac34e^x+\dfrac14e^{-x}+\dfrac{x}{2}e^x}
\else \answerarea[10] \fi%

\textbf{8.} 计算二重积分
\[
I=\int_0^1\mathrm{d}x\int_{x^2}^{1}x e^{y^2}\,\mathrm{d}y.
\]

\ifshowsolutions%
\boxed{\textbf{解}}

\textbf{思路：} 先对 $y$ 积分时 $e^{y^2}$ 没有初等原函数，需要交换积分次序；交换前把积分区域用不等式写清楚。

积分区域为
\[
D=\{(x,y)\mid 0\le x\le 1,\ x^2\le y\le 1\}.
\]
它由抛物线 $y=x^2$、直线 $y=1$ 和 $y$ 轴围成。换序时先看 $y$：$0\le y\le 1$，对固定的 $y$，$x$ 从 $0$ 变到 $\sqrt y$。因此
\[
D=\{(x,y)\mid 0\le y\le 1,\ 0\le x\le\sqrt y\}.
\]
于是
\[
I=\int_0^1\mathrm{d}y\int_0^{\sqrt y}x e^{y^2}\,\mathrm{d}x
=\int_0^1 e^{y^2}\left[\frac{x^2}{2}\right]_0^{\sqrt y}\mathrm{d}y
=\frac12\int_0^1 y e^{y^2}\,\mathrm{d}y.
\]
令 $u=y^2$，$\mathrm{d}u=2y\,\mathrm{d}y$，得
\[
I=\frac12\cdot\frac12\int_0^1 e^u\,\mathrm{d}u
=\frac14(e-1).
\]
\boxed{\text{答案：}\dfrac{e-1}{4}}
\else \answerarea[8] \fi%

\textbf{9.} 计算曲线积分
\[
I=\int_{(0,0)}^{(\pi/2,1)}
\bigl(2xy^3-y^2\cos x\bigr)\mathrm{d}x
+\bigl(3x^2y^2-2y\sin x\bigr)\mathrm{d}y,
\]
并说明该积分与路径无关。

\ifshowsolutions%
\boxed{\textbf{解}}

\textbf{思路：} 先验证 $\dfrac{\partial P}{\partial y}=\dfrac{\partial Q}{\partial x}$，在单连通区域上积分与路径无关；再求势函数 $u$，用端点值作差。

记
\[
P=2xy^3-y^2\cos x,\qquad Q=3x^2y^2-2y\sin x.
\]
则
\[
\frac{\partial P}{\partial y}=6xy^2-2y\cos x,\qquad
\frac{\partial Q}{\partial x}=6xy^2-2y\cos x.
\]
两者在整个平面（单连通区域）上相等且连续，故积分与路径无关，存在势函数 $u(x,y)$ 使
\[
u_x=P,\qquad u_y=Q.
\]
由 $u_x=2xy^3-y^2\cos x$ 对 $x$ 积分：
\[
u=x^2y^3-y^2\sin x+\varphi(y).
\]
再由
\[
u_y=3x^2y^2-2y\sin x+\varphi'(y)=Q
\]
得 $\varphi'(y)=0$，可取 $\varphi(y)=0$。因此
\[
u(x,y)=x^2y^3-y^2\sin x.
\]
于是
\[
I=u\left(\frac{\pi}{2},1\right)-u(0,0)
=\left(\frac{\pi^2}{4}\cdot1-1\cdot\sin\frac{\pi}{2}\right)-0
=\frac{\pi^2}{4}-1.
\]
也可沿折线验证：先沿 $y=0$ 从 $(0,0)$ 到 $(\pi/2,0)$，此时 $P=0$，积分为 $0$；再沿 $x=\pi/2$ 从 $y=0$ 到 $y=1$，积分为
\[
\int_0^1\left(\frac{3\pi^2}{4}y^2-2y\right)\mathrm{d}y
=\frac{\pi^2}{4}-1,
\]
与势函数法一致。
\boxed{\text{答案：}\dfrac{\pi^2}{4}-1}
\else \answerarea[10] \fi%

\subsection{真题风格与综合训练}

\textbf{10.} 设正项数列 $\{a_n\}$ 单调递减，且级数 $\displaystyle\sum_{n=1}^{\infty}a_n$ 收敛。
（1）证明：$\displaystyle\lim_{n\to\infty}na_n=0$；
（2）举例说明：若去掉“单调递减”这一条件，结论未必成立。

\ifshowsolutions%
\boxed{\textbf{解}}

\textbf{思路：} 利用级数收敛的柯西准则，把“尾段和”控制得任意小；再由单调性，尾段中每一项都不小于最后一项，从而用尾段和压住 $na_n$。奇偶下标分别处理。

（1）设 $S_n=\displaystyle\sum_{k=1}^n a_k$。由于 $\sum a_n$ 收敛，$\{S_n\}$ 是柯西列。对任意 $\varepsilon>0$，存在 $N$，使得当 $m\ge N$ 时，
\[
0\le S_{2m}-S_m=\sum_{k=m+1}^{2m}a_k<\varepsilon.
\]
因为 $\{a_n\}$ 单调递减，故
\[
a_{m+1}\ge a_{m+2}\ge\cdots\ge a_{2m},
\]
于是
\[
m\,a_{2m}\le\sum_{k=m+1}^{2m}a_k<\varepsilon,
\]
从而
\[
2m\,a_{2m}<2\varepsilon. \tag{1}
\]
对于奇数下标，
\[
(2m+1)a_{2m+1}\le(2m+1)a_{2m}
=2m\,a_{2m}+a_{2m}
<2\varepsilon+a_{2m}.
\]
又因 $\sum a_n$ 收敛，必有 $a_n\to0$，所以当 $m$ 充分大时 $a_{2m}<\varepsilon$，从而
\[
(2m+1)a_{2m+1}<3\varepsilon. \tag{2}
\]
由 (1)(2)，当 $n$ 充分大时 $0\le na_n<3\varepsilon$。由 $\varepsilon$ 的任意性，
\[
\lim_{n\to\infty}na_n=0.
\]

（2）取
\[
a_n=\begin{cases}
\dfrac{1}{k^2}, & n=k^2,\ k=1,2,\cdots,\\[4pt]
2^{-n}, & \text{其他 }n.
\end{cases}
\]
则 $a_n>0$，且
\[
\sum_{n=1}^{\infty}a_n
\le\sum_{k=1}^{\infty}\frac{1}{k^2}+\sum_{n=1}^{\infty}2^{-n}<+\infty,
\]
故级数收敛，并且 $a_n\to0$。但当 $n=k^2$ 时，
\[
n a_n=k^2\cdot\frac{1}{k^2}=1,
\]
故 $\lim\limits_{n\to\infty}na_n$ 不存在，更不等于 $0$。这个数列不是单调递减的，说明（1）中单调性的条件不能去掉。
\boxed{\text{答案：}(1)\ 0;\ (2)\ \text{反例见上}}
\else \answerarea[10] \fi%

\textbf{11.} 计算曲面积分
\[
I=\oiint_{\Sigma}x\,\mathrm{d}y\mathrm{d}z+y\,\mathrm{d}z\mathrm{d}x+z^2\,\mathrm{d}x\mathrm{d}y,
\]
其中 $\Sigma$ 为抛物面 $z=x^2+y^2$（$0\le z\le1$）与平面 $z=1$ 所围立体 $\Omega$ 的边界曲面，取外侧。

\ifshowsolutions%
\boxed{\textbf{解}}

\textbf{思路：} 曲面是封闭的且取外侧，直接用高斯公式化为三重积分；立体在 $xy$ 平面上的截面随 $z$ 变化，用柱坐标（先 $r$ 后 $z$）最方便。

令
\[
P=x,\qquad Q=y,\qquad R=z^2.
\]
则
\[
\frac{\partial P}{\partial x}+\frac{\partial Q}{\partial y}+\frac{\partial R}{\partial z}
=1+1+2z=2+2z.
\]
由高斯公式，
\[
I=\iiint_{\Omega}(2+2z)\,\mathrm{d}V.
\]
立体 $\Omega$ 可表示为
\[
0\le z\le1,\qquad 0\le r\le\sqrt z,\qquad 0\le\theta\le2\pi.
\]
于是
\[
\begin{aligned}
I&=\int_0^{2\pi}\mathrm{d}\theta\int_0^1\mathrm{d}z\int_0^{\sqrt z}(2+2z)r\,\mathrm{d}r\\
&=2\pi\int_0^1(2+2z)\cdot\frac{z}{2}\,\mathrm{d}z\\
&=2\pi\int_0^1(z+z^2)\,\mathrm{d}z
=2\pi\left(\frac12+\frac13\right)
=\frac{5\pi}{3}.
\end{aligned}
\]
\boxed{\text{答案：}\dfrac{5\pi}{3}}
\else \answerarea[10] \fi%

\textbf{12.} 设 $f(x)$ 在 $[0,1]$ 上二阶连续可导，$f(0)=f(1)=0$，且 $f$ 在 $[0,1]$ 上的最小值为 $-1$。证明：
\[
\max_{0\le x\le1}|f''(x)|\ge 8.
\]

\ifshowsolutions%
\boxed{\textbf{解}}

\textbf{思路：} 最小值在内部某点 $c$ 取得，故 $f'(c)=0$。分别把 $f(0)$、$f(1)$ 在 $c$ 处用带拉格朗日余项的泰勒公式展开，得到两个关于 $f''$ 的等式，再按 $c\le\frac12$ 与 $c\ge\frac12$ 分类估计。

设 $f$ 在 $c\in[0,1]$ 处取得最小值 $-1$。由于 $f(0)=f(1)=0$，最小值不可能在端点取得，故 $c\in(0,1)$，且由费马定理
\[
f'(c)=0.
\]

将 $f(0)$ 在 $x=c$ 处展开：
\[
f(0)=f(c)+f'(c)(0-c)+\frac{f''(\xi_1)}{2}(0-c)^2,
\]
其中 $\xi_1\in(0,c)$。代入 $f(0)=0,\ f(c)=-1,\ f'(c)=0$，得
\[
0=-1+\frac{f''(\xi_1)}{2}c^2,
\]
即
\[
f''(\xi_1)=\frac{2}{c^2}. \tag{1}
\]
同理，将 $f(1)$ 在 $x=c$ 处展开：
\[
0=f(1)=-1+\frac{f''(\xi_2)}{2}(1-c)^2,
\]
其中 $\xi_2\in(c,1)$，故
\[
f''(\xi_2)=\frac{2}{(1-c)^2}. \tag{2}
\]

若 $0<c\le\dfrac12$，则 $\dfrac1{c^2}\ge4$，由 (1) 得
\[
|f''(\xi_1)|=\frac{2}{c^2}\ge8.
\]
若 $\dfrac12\le c<1$，则 $1-c\le\dfrac12$，由 (2) 得
\[
|f''(\xi_2)|=\frac{2}{(1-c)^2}\ge8.
\]
无论哪种情况，都存在 $\xi\in(0,1)$ 使 $|f''(\xi)|\ge8$，所以
\[
\max_{0\le x\le1}|f''(x)|\ge8.
\]

等号可以取到：$f(x)=4x(x-1)$ 满足 $f(0)=f(1)=0$，在 $x=\frac12$ 处取最小值 $-1$，且 $f''(x)=8$，故常数 $8$ 是最佳的。
\boxed{\text{答案：}\text{结论成立，常数 }8\text{ 不能再改进}}
\else \answerarea[10] \fi%
